\section{서론}

배틀로얄 경기의 최종 순위는 플레이 가능 영역이 축소되면서 변하는
생존 위험이 누적된 결과이다. \textit{PUBG: Battlegrounds}에서
플레이어는 낙하하고, 아이템을 획득하고, 교전하며, 점점 줄어드는
전장에서 마지막까지 생존해야 한다. 블루존 바깥에 오래 머무르면
체력이 감소하므로 플레이어는 계속 이동 압력을 받는다
~\cite{pubgGameInfo}. 경기 중 안전 구역 갱신과 블루존 이동은
텔레메트리에서 관측 가능한 존 변화 구간을 형성한다. 본 연구는
각 구간에 인덱스를 부여하여 \textit{분석 페이즈}로 사용함으로써, 공간 수축
과정의 비교 가능한 시점에서 생존 위험을 측정한다. 이 표현에서는 같은
피해량, 이동 거리, 적과의 조우도 현재의 존 압력과 안전 구역 기하에
따라 다른 위험을 의미한다. 따라서 경기 종료 후 집계 통계만으로는 팀이
언제 취약해졌고 어떤 공간 조건이 그 위험을 만들었는지 설명하기 어렵다.

본 연구는 공식 PUBG 텔레메트리를 팀 관점에서 분석한다. 개발자 API는
시간에 따른 플레이어 위치, 교전, 존 상태 변화 이벤트, 탈락을
기록한다~\cite{pubgDeveloperPortal,pubgTelemetryEvents,pubgTelemetryObjects}.
우리는 플레이어 기록을 팀 궤적, 존 노출, 팀 간 관계로 집계하고,
마지막 생존 팀원이 탈락한 시점을 팀 탈락 시점으로 정의한다. 핵심 예측
문제는 인덱싱된 분석 구간 \(k\) 종료 시점까지 관측 가능한 정보만으로
해당 팀이 다음 구간 \(k+1\) 동안 탈락하는지를 예측하는 것이다. 이를
이산시간 event-history 관점의 \textit{다음 페이즈 탈락 위험}으로 정의한다
~\cite{allison1982discrete,singer1993discrete}.

본 논문은 두 질문을 다룬다. \textbf{RQ1}: 공식 PUBG 텔레메트리를
페이즈 정렬 STDB로 재구성하여 팀 궤적, 존 노출, 교전 이벤트, 동적
팀 관계를 특징과 그래프로 도출할 수 있는가? \textbf{RQ2}: 이 표현
위에서 특징 기반 모델은 다음 페이즈 탈락 위험을 얼마나 잘 예측하며,
동적 그래프는 특징만으로 설명되지 않는 추가 신호를 제공하는가?
기여는 세 가지이다. 첫째, 팀 단위 페이즈 수준 생존 위험 예측 문제를
정의한다. 둘째, 21,824개 경기의 페이즈 정렬 STDB 표현을 구축한다.
셋째, 표현 자체의 이득과 그래프 관계의 추가 이득을 실험적으로
분리한다.
